\section{The entropy inequality}\label{sec:03-entropy}

We use two cone-valued maps of a Gaussian to bound the dual Laplace
integrals. Change of variables bounds these integrals through the maps' expected
log Jacobians and linear costs. Comparison with the one-dimensional
Gaussian model will put the remaining scalar terms into the form needed
for the quadratic argument.
We first prove the log integrability required by change of variables
and Jensen's inequality, and then carry out that argument with
smoothed maps.

\subsection{Projection maps and their Jacobians}
In the orthant example of Section~\ref{sec:02-projections}, the
coordinate maps $-\log p(-G_i)$ and $-\log p(G_i)$ take values in the
positive half-line. Since $p(G_i)$ is uniform on $(0,1)$, each has the
standard exponential distribution. The following weights express
these maps as integrals of the biased projections and extend the
construction to arbitrary cones.

Define the positive scalar weights
\[
\begin{gathered}
c_X(z)=\frac{\phi(z)}{p(z)},\qquad j_X(z)=z+c_X(z),\qquad
W_X(z)=-c_X'(z)=j_X(z)c_X(z),\\
c_Y(z)=c_X(-z),\qquad W_Y(z)=W_X(-z)=c_Y'(z),
\end{gathered}
\]
and the random positive semidefinite matrices
\[
\mathcal J_X=\int_{\mathbb R}W_X(z)P_z\,dz,\qquad
\mathcal J_Y=\int_{\mathbb R}W_Y(z)(I-P_z)\,dz.
\]
The weights are polynomially bounded; \(W_X\) decays at a Gaussian rate at the positive end, and \(W_Y\) at the negative end. Thus these matrices have finite expected trace.


Consider the formal cone-valued maps
\[
 \mathcal T_X(G)=\int_{\mathbb R}W_X(z)X_z\,dz,
 \qquad
 \mathcal T_Y(G)=\int_{\mathbb R}W_Y(z)Y_z\,dz.
\]
Their formal Jacobian factors are
$\mathcal J_X\Sigma^{1/2}$ and $-\mathcal J_Y\Sigma^{1/2}$. In the orthant model, integration by
parts gives exactly the coordinate maps above. To justify change
of variables without imposing regularity on the cone boundary, we
will use truncated, Gaussian-smoothed versions of these maps.
The first task is to control the log determinants of
$\mathcal J_X$ and $\mathcal J_Y$.

Define the symmetric nonnegative kernel
\[
 \operatorname{nt}(x,z)
 =\tau\bigl(P_{\min(x,z)}\circ(I-P_{\max(x,z)})\bigr),
 \qquad
 N=\iint\operatorname{nt}(x,z)\,dx\,dz.
\]
Its nonnegativity follows because the trace of a product of two
positive semidefinite matrices is nonnegative. It is integrable
against any polynomial weight: on $x\le z$, split into
$x\le z\le0$, $x\le0\le z$, and $0\le x\le z$.
Lemma~\ref{lem:projection-tails} gives Gaussian decay of $P_x$
at the negative end and of $I-P_z$ at the positive end.
On the mixed region H\"older's inequality gives decay in both
variables; on each same-sign region the inner interval has
polynomial length. The other ordered half-plane follows by symmetry.

\subsection{The log-Jacobian estimate}
\begin{lemma}[Log-Jacobian gain]\label{lem:log-jacobian}
The matrices \(\mathcal J_X,\mathcal J_Y\) are positive definite almost surely, their log determinants are integrable, and
\[
\tau(\log\mathcal J_X+\log\mathcal J_Y)
\ge-\gamma v-\tau H_v+\tfrac18N.
\tag{9}\label{eq:11}
\]
\end{lemma}
\begin{proof}
We first establish integrability and an exact formula for the contribution of \(\mathcal J_X\). Set
\[
J(z)=\int_{-\infty}^zW_X(l)P_l\,dl,\qquad
S_z=J(z)/c_X(z),\qquad F_X(z)=\tau\log(c_X(z)I+J(z)).
\]
For finite \(z\) the matrix inside the logarithm is bounded below by \(c_X(z)I\), and its expected trace is finite. Hence its log determinant is integrable. Its derivative is \(-W_X(z)(I-P_z)\), so the matrix and \(F_X\) decrease with \(z\). The trace differentiation formula gives, almost everywhere,
\[
F_X'(z)=-j_X(z)(1-h_1(z))+\beta_X(z),\qquad
h_1(z)=\tau P_z,
\]
where
\[
\beta_X(z)=j_X(z)\tau\bigl[S_z(I+S_z)^{-1}(I-P_z)\bigr]\ge0.
\]
The derivative is locally bounded uniformly in the Gaussian input after normalization, which justifies differentiation under expectation on compact intervals. The nonnegativity uses positivity under the trace and requires no commutation.

Compare \(F_X\) with
\[
F_0(z)=\int\log c_X(\min(z,x))\,d\gamma(x),\qquad
F_0'=-j_X(1-p).
\]
Both functions equal \(\log c_X(z)+o(1)\) at negative infinity. For \(F_X\), this follows from
\(0\le\tau\log(I+S_z)\le\tau S_z=o(1)\), using the layer tails; for \(F_0\), it follows from the Gaussian tail and the polynomial growth of \(\log c_X\). Thus
\[
(F_X-F_0)'=j_X(h_1-p)+\beta_X,
\]
and the signed first term is absolutely integrable. Integrating from a finite lower endpoint and letting it tend to negative infinity proves that \(\beta_X\) is integrable on every left half-line. At positive infinity \(F_0\) tends to the finite number \(\gamma\log c_X\). If the remaining integral of \(\beta_X\) were infinite, the last identity would force \(F_X\) to tend to positive infinity, contrary to its decrease. The same identity bounds \(F_X\) below. Consequently the limiting expectation is finite and
\[
\lim_{z\to\infty}F_X(z)
=\gamma\log c_X+\int j_X(h_1-p)\,dz+\int\beta_X(z)\,dz.
\]
The matrices decrease to \(\mathcal J_X\). Subtracting the integrable log determinant at any finite layer permits monotone convergence to their limit, even if a zero eigenvalue is initially allowed there. The finite limiting expectation excludes that possibility almost surely. The finite expected trace controls the positive part of the limiting log determinant, so its negative part is integrable as well. We have therefore proved
\[
\tau\log\mathcal J_X
=\gamma\log c_X+\int j_X(h_1-p)\,dz+\int\beta_X(z)\,dz.
\]
Reflection, replacing \(P_z\) by \(I-P_{-z}\), gives the corresponding identity for \(\mathcal J_Y\).

It remains to estimate the nonnegative surplus. For \(l\le z\), set
\(D_{l,z}=c_X(z)^{-1}\int_l^zW_X(y)P_y\,dy\). Then
\[
0\le D_{l,z}\le S_z,\qquad
D_{l,z}\le\bigl(c_X(l)/c_X(z)-1\bigr)I.
\]
The matrix function \(D\mapsto D(I+D)^{-1}=I-(I+D)^{-1}\) is order preserving. Applying this fact and then scalar spectral calculus to \(D_{l,z}\) yields
\[
S_z(I+S_z)^{-1}
\ge D_{l,z}(I+D_{l,z})^{-1}
\ge\frac1{c_X(l)}\int_l^zW_X(y)P_y\,dy.
\]

We will average this inequality over \(l\). First, \(W_X\) is decreasing. To prove this, let \(g_1\) be standard normal and condition on \(g_1\le z\). The nonnegative random variable \(z-g_1\) has mean \(j_X(z)\) and variance
\(j_X'(z)=1-c_X(z)j_X(z)>0\). Its survival function
\[
H_*(s)=p(z-s)/p(z),\qquad s\ge0,
\]
is log-concave, because \((\log p)''=-j_Xc_X\). Since \(H_*(0)=1\), concavity of its logarithm implies
\(H_*(s+t)\le H_*(s)H_*(t)\). Integrating over the positive quadrant gives
\[
\tfrac12\mathbb E[(z-g_1)^2\mid g_1\le z]
=\int_0^\infty\!\int_0^\infty H_*(s+t)\,ds\,dt
\le j_X(z)^2.
\]
Thus \(j_X'\le j_X^2\), and \(W_X'=c_X(j_X'-j_X^2)\le0\).

For \(z\ge0\), choose the positive measure \(\omega_z\) on \([-z,z]\) specified by
\[
\int_{[-z,y]}\frac{d\omega_z(l)}{c_X(l)}
=\frac1{4j_X(z)W_X(y)},\qquad -z\le y\le z.
\]
Its existence follows from the monotonicity of \(1/W_X\), with the indicated initial atom at \(-z\). Integration by parts gives
\[
\omega_z([-z,z])
=\frac1{4j_X(z)^2}+\frac{z}{2j_X(z)}
\le\frac\pi8+\frac12<1,
\]
since \(j_X\) is increasing, \(j_X(0)=\sqrt{2/\pi}\), and \(j_X(z)\ge z\). Integrating the matrix inequality against this subprobability measure and using positivity gives
\[
\beta_X(z)\ge\frac14\int_{-z}^z\operatorname{nt}(y,z)\,dy\qquad(z\ge0).
\]
The reflected argument gives the corresponding estimate on the other half of the ordered region \(y\le z\): the two regions are distinguished by \(y+z\ge0\) and \(y+z\le0\). Their common boundary has measure zero. Since \(N\) is the integral of a symmetric kernel over the whole plane, the two surplus integrals sum to at least \(N/8\).

Finally, the two scalar comparison terms sum to \(-\gamma v\), because \(c_Xc_Y=1/(XY)\). The linear terms sum to
\[
\int\bigl(j_X(z)-j_X(-z)\bigr)(h_1(z)-p(z))\,dz=-\tau H_v,
\]
using \(j_X(z)-j_X(-z)=v'(z)\). This proves Equation~\eqref{eq:11}.
\end{proof}

\subsection{Cone-valued maps and change of variables}
\begin{lemma}[Entropy from the projection maps]\label{lem:cone-entropy-maps}
With \(C_Y=\int W_YB=\int c_Yr_1\),
\[
\frac1m\log\bigl(\chi_C(V)\chi_D(U)\bigr)
\ge\log(2\pi e)+\tau\log\Sigma
+\tau(\log\mathcal J_X+\log\mathcal J_Y)-1-C_Y.
\]
\end{lemma}
\begin{proof}
Let \(Z'\) be an independent copy of \(Z\), fix \(0<\zeta<1\), and put
\(\widetilde Z=\zeta Z+\sqrt{1-\zeta^2}Z'\). The law of \(\widetilde Z\) is the same as that of \(Z\). Temporarily write \(X_z(h),Y_z(h),P_z(h)\) for the fields evaluated with \(Z=h\). For integers \(j\ge1\), define maps of \(G\) by
\[
\begin{aligned}
\mathcal T_{X,j}(G)&=\int_{-j}^jW_X(z)\mathbb E_{Z'}X_z(\widetilde Z)\,dz,\\
\mathcal T_{Y,j}(G)&=\int_{-j}^jW_Y(z)\mathbb E_{Z'}Y_z(\widetilde Z)\,dz.
\end{aligned}
\]
Their Jacobians are \(J_{X,j}\zeta\Sigma^{1/2}\) and \(-J_{Y,j}\zeta\Sigma^{1/2}\), where
\[
\begin{aligned}
J_{X,j}&=\int_{-j}^jW_X(z)\mathbb E_{Z'}P_z(\widetilde Z)\,dz,\\
J_{Y,j}&=\int_{-j}^jW_Y(z)\mathbb E_{Z'}(I-P_z(\widetilde Z))\,dz.
\end{aligned}
\]
Gaussian smoothing makes the maps smooth. Both factors are positive definite, because the conditional Gaussian sees an open set where the relevant projection derivative is the identity. Viewed as functions of \(Z\), the two maps have symmetric positive and negative definite derivatives, respectively. Integration along line segments proves strict monotonicity up to sign, hence injectivity. Their values belong to \(C\) and \(D\); since their derivatives are nonsingular, their images are open and lie in the interiors of these cones. The inverse function theorem gives a smooth inverse on each image.

For fixed \(j,\zeta\), the log determinants of the factors are integrable. Indeed, restrict the layer integral to a fixed compact subinterval of \([-1,1]\). Its biases are bounded. A fixed ball in \(\operatorname{int}C\), or in \(-\operatorname{int}D\), has conditional Gaussian probability at least \(c\exp[-C(1+|Z|)^2]\) after the relevant translations. Thus each factor is bounded below by this scalar times \(I\). The upper bound is deterministic, since \(0\le P_z\le I\) and the layer interval is finite. The map values also have finite first moments by the Lipschitz bounds.

For completeness, let \(\rho\) be the standard Gaussian density on \(\mathbb R^m\). Change of variables over the image of \(\mathcal T_{X,j}\), followed by Jensen's inequality, gives
\[
\begin{aligned}
\log\chi_C(V)
&\ge\log\mathbb E\frac{e^{-V\cdot\mathcal T_{X,j}(G)}
 |\det D_G\mathcal T_{X,j}(G)|}{\rho(G)}\\
&\ge\tfrac m2\log(2\pi e)
+\mathbb E\log|\det D_G\mathcal T_{X,j}|-\mathbb E[V\cdot\mathcal T_{X,j}].
\end{aligned}
\]
The integrability just proved justifies this step. The same argument applies to \(\mathcal T_{Y,j}\) and \(\chi_D(U)\). It uses only containment of the images in the cones, not surjectivity.

As \(j\to\infty\), the factors increase in matrix order to
\(\mathbb E_{Z'}\mathcal J_X(\widetilde Z)\) and
\(\mathbb E_{Z'}\mathcal J_Y(\widetilde Z)\).
Their expected traces are finite by equality of the marginal laws. Subtracting the integrable log determinant for \(j=1\) makes monotone convergence applicable; the positive parts are controlled by the expected traces. Conditional concavity of log determinant \cite{BoydVandenberghe2004} then gives
\[
\mathbb E\log\det\mathbb E_{Z'}\mathcal J_X(\widetilde Z)
\ge\mathbb E\log\det\mathcal J_X(Z),
\]
and similarly for \(Y\). Lemma~\ref{lem:log-jacobian} supplies the two-sided integrability needed for this conditional Jensen inequality.

After dividing the sum of the two bounds by \(m\), the two right Jacobian factors contribute \(2\log\zeta+\tau\log\Sigma\). By the marginal law of \(\widetilde Z\), Equation~\eqref{eq:5}, and \(U\cdot V=m\), the cost terms converge to \(\int W_Xa\) and \(\int W_YB\). The layer tails justify both limits. Integration by parts gives
\[
\int W_Xa=\int c_Xp=\int\phi=1,\qquad
\int W_YB=\int c_Yr_1=C_Y.
\]
The boundary products vanish by the same tails. Letting \(\zeta\uparrow1\) removes the term \(2\log\zeta\) and completes the proof. No limiting change of variables for the unsmoothed maps is needed.
\end{proof}

\subsection{Identification of the scalar terms}
Combining Lemmas~\ref{lem:log-jacobian} and~\ref{lem:cone-entropy-maps} gives
\[
\frac1m\log\bigl(\chi_C(V)\chi_D(U)\bigr)
\ge1+\tau\log\Sigma-\tau H_v-C_Y+\tfrac18N.
\]
Here \(\gamma v=\log(2\pi)-1\): under \(\gamma\), \(p\) is uniform on \((0,1)\), so each of \(\gamma\log p\) and \(\gamma\log(1-p)\) equals \(-1\), whereas \(\gamma(-2\log\phi)=\log(2\pi)+1\).

We now compare $B,r_1,h$, and $\mu$ with their Gaussian references.
This will express the scalar cost in terms of the profile $d$.

The Gaussian reference functions are
\[
B_0(z)=a(-z),\qquad q_0(z)=1-p(z).
\]
Put
\[
\Delta B=B-s_0B_0,\quad \Delta r=r_1-s_0q_0,\quad
h_\delta=h-s_0p,\quad \Delta_\mu=\mu-s_0\gamma.
\]
The identity \(pB_0+aq_0=\phi\) shows that \(\Delta_\mu\) has density \(p\Delta B+a\Delta r=h_\delta'\); also \((\Delta B)'=-\Delta r\). As for \(\gamma\), we use measures as integration functionals. For a smooth scalar test \(j\), define
\[
\delta[j]=-
\int_{\mathbb R}\bigl[2p\Delta B\,j+(p\Delta r+h_\delta)(2zj-j')\bigr]\,dz.
\tag{10}\label{eq:9}
\]
This integral is absolutely convergent. In particular \(\Delta B\) grows at most linearly at the negative end, while \(p\Delta B\) and \(h_\delta\) have Gaussian tails; the weighted integrability of \(p\Delta r\) follows from Lemma~\ref{lem:projection-tails}.

\begin{proposition}[Entropy inequality]\label{prop:entropy}
For the cone data and the fixed projection field constructed in Section~\ref{sec:02-projections},
\[
\frac1m\log\bigl(\chi_C(V)\chi_D(U)\bigr)
\ge \tau(\log\Sigma-T)+\tau_T H_v-\delta[d]+\tfrac18N.
\tag{11}\label{eq:10}
\]
\end{proposition}

\begin{proof}[Proof of Proposition~\ref{prop:entropy}]
Integration by parts gives
\[
\tau_\Sigma H_v=-\int h_\delta v'\,dz=\Delta_\mu v.
\]
We will prove \(C_Y-s_0+\Delta_\mu v=\delta[d]\). First,
\(\int p\Delta B=\int a\Delta r\), by integrating the derivative of \(a\Delta B\); their sum is the zero total mass of \(\Delta_\mu\). Thus both integrals vanish. Since \(\int c_Yq_0=1\) and \(c_Y=2a-pv'\),
\[
C_Y-s_0+\Delta_\mu v=-\int(p\Delta r+h_\delta)v'\,dz.
\]

For any smooth polynomial-growth test \(w_0\), we have
\[
\int\bigl[p\Delta B(2+\mathcal N)w_0
 +(p\Delta r+h_\delta)w_0'\bigr]\,dz=0.
\tag{12}\label{eq:12}
\]
Indeed, integrate the \(w_0''\) term once by parts and use
\((\Delta B)'=-\Delta r\), \(h_\delta'=p\Delta B+a\Delta r\), and \(a-\phi=zp\). The expression becomes
\[
\int\bigl[2p\Delta B\,w_0+(a\Delta B+h_\delta)w_0'\bigr]\,dz,
\]
which is zero because \((a\Delta B+h_\delta)'=2p\Delta B\). The boundary products vanish: \(\Delta B=O(1+|z|)\) at the negative end, where \(a,p\) have Gaussian decay, and \(B,h-s_0\) have Gaussian decay at the positive end. The densities involving \(\Delta r\) have all needed polynomial moments by Lemma~\ref{lem:projection-tails}.

Apply Equation~\eqref{eq:12} to \(w_0=g\). Equation~\eqref{eq:2} gives
\[
\int\bigl[p\Delta B\,d+(p\Delta r+h_\delta)g'\bigr]\,dz=0.
\]
Substituting \(2zd-d'=v'+2g'\) in the definition of \(\delta[d]\) now proves the desired scalar identity. Finally,
\[
\tau H_v=\Delta_\mu v-\tau_T H_v,\qquad s_0=1+\tau T.
\]
Inserting these identities in the entropy bound yields exactly Equation~\eqref{eq:10}.
\end{proof}

Define
\[
 \mathcal E=\delta[d]-\tau_T H_v+\tau(T-\log\Sigma)-\frac18N.
\]
Equation~\eqref{eq:10} bounds the normalized logarithm of the product
of Laplace integrals below by $-\mathcal E$. It therefore remains to
prove $\mathcal E\le0$.
